% GENERATED FILE. Edit canonical lessons, then run npm run generate:books.
% canonical-source-hash: fnv1a64:407f15e56008efa8
% canonical-lessons: TA-C410-patacari, TA-C410-valippokkar, TA-C410-palakai, TA-C410-ulavu, TA-C410-tantu

\chapter{Walking Through Town}
\label{ch:ta-walking-through-town}
\hlchaptermodality{410}

\begin{chapteropening}
\textbf{By the end of this chapter:} \emph{I can talk about pedestrians, passers-by and signboards, and strolling and jumping over.}

Complete the last lesson of chapter 410: I can talk about pedestrians, passers-by and signboards, and strolling and jumping over.
\end{chapteropening}

\section[\texorpdfstring{\ta{பாதசாரி} (pātacāri)}{பாதசாரி (pātacāri)}]{\ta{பாதசாரி} (pātacāri) — a pedestrian}

\label{lesson:TA-C410-patacari}



\begin{quote}
Before the new one: say the Tamil for a shortcut, then the Tamil for to get stuck.
\end{quote}

\subsection*{You'll want to know: \ta{பாதசாரி}}
\textbf{\ta{பாதசாரி}} — \emph{pātacāri} — \textquotedblleft{}a pedestrian\textquotedblright{}.

\begin{quote}
\ta{பாதசாரிகள்} \ta{நடைபாதையில்} \ta{நடக்க} \ta{வேண்டும்}.
\end{quote}

\begin{quote}
\emph{pātacārikaḷ naṭaipātaiyil naṭakka vēṇṭum.}
\end{quote}

\begin{quote}
Pedestrians should walk on the footpath.
\end{quote}

The plural is \ta{பாதசாரிகள்}.

\begin{cousinweb}
\textbf{\ta{பாதசாரி}} (\emph{pātacāri}) is borrowed from Sanskrit \emph{pādacārin}, \textquotedblleft{}going on foot, a pedestrian\textquotedblright{} (Monier-Williams), from \emph{pāda}, \textquotedblleft{}foot\textquotedblright{}.
\end{cousinweb}

\begin{grammarlens}[title={What you've built}]
One more word for getting around.
\end{grammarlens}

\subsection*{Your turn}
\begin{itemize}
\raggedright
  \item \emph{Say it:} \emph{pātacāri}
  \item \emph{Say it:} \emph{pātacāri}, once more
  \item \emph{Say it:} \emph{pātacārikaḷ naṭaipātaiyil naṭakka vēṇṭum.} — Pedestrians should walk on the footpath.
  \item \emph{Recall:} say the Tamil for a shortcut, then the Tamil for to get stuck, then say \emph{pātacāri} again
\end{itemize}

\begin{tcolorbox}[breakable,colback=teal!4,colframe=teal!35!black,title={Before you move on}]
What does \ta{பாதசாரி} mean? (\textquotedblleft{}A pedestrian\textquotedblright{}.) Say it once more.
\end{tcolorbox}
\section[\texorpdfstring{\ta{வழிப்போக்கர்} (vaḻippōkkar)}{வழிப்போக்கர் (vaḻippōkkar)}]{\ta{வழிப்போக்கர்} (vaḻippōkkar) — a passer-by; a wayfarer}

\label{lesson:TA-C410-valippokkar}



\begin{quote}
Before the new one: say the Tamil for to get stuck, then the Tamil for a pedestrian.
\end{quote}

\subsection*{You'll want to know: \ta{வழிப்போக்கர்}}
\textbf{\ta{வழிப்போக்கர்}} — \emph{vaḻippōkkar} — \textquotedblleft{}a passer-by; a wayfarer\textquotedblright{}.

\begin{quote}
\ta{ஒரு} \ta{வழிப்போக்கர்} \ta{எனக்கு} \ta{வழி} \ta{காட்டினார்}.
\end{quote}

\begin{quote}
\emph{oru vaḻippōkkar eṉakku vaḻi kāṭṭiṉār.}
\end{quote}

\begin{quote}
A passer-by showed me the way.
\end{quote}

\ta{வழி}, way, + \ta{போக்கு}, going, + -\ta{அர்}, a person.

\begin{cousinweb}
\textbf{\ta{வழிப்போக்கர்}} (\emph{vaḻippōkkar}) is built from \ta{வழி}, \textquotedblleft{}a way\textquotedblright{} (DEDR 5297), \ta{போக்கு}, \textquotedblleft{}going\textquotedblright{} (DEDR 4572, the family of \ta{போ}), and the personal ending \emph{-ar}.
\end{cousinweb}

\begin{grammarlens}[title={What you've built}]
One more word for getting around.
\end{grammarlens}

\subsection*{Your turn}
\begin{itemize}
\raggedright
  \item \emph{Say it:} \emph{vaḻippōkkar}
  \item \emph{Say it:} \emph{vaḻippōkkar}, once more
  \item \emph{Say it:} \emph{oru vaḻippōkkar eṉakku vaḻi kāṭṭiṉār.} — A passer-by showed me the way.
  \item \emph{Recall:} say the Tamil for to get stuck, then the Tamil for a pedestrian, then say \emph{vaḻippōkkar} again
\end{itemize}

\begin{tcolorbox}[breakable,colback=teal!4,colframe=teal!35!black,title={Before you move on}]
What does \ta{வழிப்போக்கர்} mean? (\textquotedblleft{}A passer-by; a wayfarer\textquotedblright{}.) Say it once more.
\end{tcolorbox}
\section[\texorpdfstring{\ta{பலகை} (palakai)}{பலகை (palakai)}]{\ta{பலகை} (palakai) — a board; a plank}

\label{lesson:TA-C410-palakai}



\begin{quote}
Before the new one: say the Tamil for a pedestrian, then the Tamil for a passer-by.
\end{quote}

\subsection*{You'll want to know: \ta{பலகை}}
\textbf{\ta{பலகை}} — \emph{palakai} — \textquotedblleft{}a board; a plank\textquotedblright{}.

\begin{quote}
\ta{தெருவின்} \ta{பெயர்} \ta{ஒரு} \ta{பலகையில்} \ta{இருக்கிறது}.
\end{quote}

\begin{quote}
\emph{teruviṉ peyar oru palakaiyil irukkiṟatu.}
\end{quote}

\begin{quote}
The name of the street is on a board.
\end{quote}

\ta{பெயர்ப்} \ta{பலகை} is a nameboard; \ta{கரும்பலகை}, a blackboard, is built on it.

\begin{cousinweb}
\textbf{\ta{பலகை}} (\emph{palakai}) is borrowed from Sanskrit \emph{phalaka}, \textquotedblleft{}a board, a plank\textquotedblright{} (CDIAL 9053).
\end{cousinweb}

\begin{grammarlens}[title={What you've built}]
One more word for getting around.
\end{grammarlens}

\subsection*{Your turn}
\begin{itemize}
\raggedright
  \item \emph{Say it:} \emph{palakai}
  \item \emph{Say it:} \emph{palakai}, once more
  \item \emph{Say it:} \emph{teruviṉ peyar oru palakaiyil irukkiṟatu.} — The name of the street is on a board.
  \item \emph{Recall:} say the Tamil for a pedestrian, then the Tamil for a passer-by, then say \emph{palakai} again
\end{itemize}

\begin{tcolorbox}[breakable,colback=teal!4,colframe=teal!35!black,title={Before you move on}]
What does \ta{பலகை} mean? (\textquotedblleft{}A board; a plank\textquotedblright{}.) Say it once more.
\end{tcolorbox}
\section[\texorpdfstring{\ta{உலாவு} (ulāvu)}{உலாவு (ulāvu)}]{\ta{உலாவு} (ulāvu) — to stroll; to wander about}

\label{lesson:TA-C410-ulavu}



\begin{quote}
Before the new one: say the Tamil for a passer-by, then the Tamil for a board.
\end{quote}

\subsection*{You'll want to know: \ta{உலாவு}}
\textbf{\ta{உலாவு}} — \emph{ulāvu} — \textquotedblleft{}to stroll; to wander about\textquotedblright{}.

\begin{quote}
\ta{மாலையில்} \ta{நாங்கள்} \ta{கடற்கரையில்} \ta{உலாவினோம்}.
\end{quote}

\begin{quote}
\emph{mālaiyil nāṅkaḷ kaṭaṟkaraiyil ulāviṉōm.}
\end{quote}

\begin{quote}
In the evening we strolled on the beach.
\end{quote}

Past (I …): \ta{உலாவினேன்}, \emph{ulāviṉēṉ}.

Where one strolls takes -\ta{இல்}: \ta{கடற்கரையில்} \ta{உலாவு}.

\begin{cousinweb}
\textbf{\ta{உலாவு}} (\emph{ulāvu}) is native Dravidian (DEDR 2693), \textquotedblleft{}to move about, to walk or ride for pleasure\textquotedblright{}.
\end{cousinweb}

\begin{grammarlens}[title={What you've built}]
One more word for getting around.
\end{grammarlens}

\subsection*{Your turn}
\begin{itemize}
\raggedright
  \item \emph{Say it:} \emph{ulāvu}
  \item \emph{Say it:} \emph{ulāvu}, once more
  \item \emph{Say it:} \emph{mālaiyil nāṅkaḷ kaṭaṟkaraiyil ulāviṉōm.} — In the evening we strolled on the beach.
  \item \emph{Recall:} say the Tamil for a passer-by, then the Tamil for a board, then say \emph{ulāvu} again
\end{itemize}

\begin{tcolorbox}[breakable,colback=teal!4,colframe=teal!35!black,title={Before you move on}]
What does \ta{உலாவு} mean? (\textquotedblleft{}To stroll; to wander about\textquotedblright{}.) Say it once more.
\end{tcolorbox}
\section[\texorpdfstring{\ta{தாண்டு} (tāṇṭu)}{தாண்டு (tāṇṭu)}]{\ta{தாண்டு} (tāṇṭu) — to cross over; to jump over}

\label{lesson:TA-C410-tantu}



\begin{quote}
Before the new one: say the Tamil for a board, then the Tamil for to stroll.
\end{quote}

\subsection*{You'll want to know: \ta{தாண்டு}}
\textbf{\ta{தாண்டு}} — \emph{tāṇṭu} — \textquotedblleft{}to cross over; to jump over\textquotedblright{}.

\begin{quote}
\ta{சிறுவன்} \ta{சிறிய} \ta{ஓடையைத்} \ta{தாண்டினான்}.
\end{quote}

\begin{quote}
\emph{ciṟuvaṉ ciṟiya ōṭaiyait tāṇṭiṉāṉ.}
\end{quote}

\begin{quote}
The boy jumped over the little stream.
\end{quote}

Past (I …): \ta{தாண்டினேன்}, \emph{tāṇṭiṉēṉ}.

What is crossed takes -\ta{ஐ}: \ta{ஓடையைத்} \ta{தாண்டு}.

\begin{cousinweb}
\textbf{\ta{தாண்டு}} (\emph{tāṇṭu}) is native Dravidian (DEDR 3158), \textquotedblleft{}to jump over, to cross, to step over\textquotedblright{}.
\end{cousinweb}

\begin{grammarlens}[title={What you've built}]
One more word for getting around.
\end{grammarlens}

\subsection*{Your turn}
\begin{itemize}
\raggedright
  \item \emph{Say it:} \emph{tāṇṭu}
  \item \emph{Say it:} \emph{tāṇṭu}, once more
  \item \emph{Say it:} \emph{ciṟuvaṉ ciṟiya ōṭaiyait tāṇṭiṉāṉ.} — The boy jumped over the little stream.
  \item \emph{Recall:} say the Tamil for a board, then the Tamil for to stroll, then say \emph{tāṇṭu} again
\end{itemize}

\begin{tcolorbox}[breakable,colback=teal!4,colframe=teal!35!black,title={Before you move on}]
What does \ta{தாண்டு} mean? (\textquotedblleft{}To cross over; to jump over\textquotedblright{}.) Say it once more.
\end{tcolorbox}
